\chapter{Male-pattern baldness}

\section*{Definition and Overview}
Male-pattern baldness, scientifically termed androgenetic alopecia, is the most common form of hair loss in men. It is characterised by a distinctive, progressive thinning of the hair on the scalp, typically presenting as a receding hairline at the temples and thinning at the crown (the vertex), which may eventually merge to leave a horseshoe-shaped pattern of hair. This condition is driven by a combination of genetic factors and male sex hormones, particularly dihydrotestosterone (DHT), which causes hair follicles to shrink over time.

\section*{Epidemiology}
Androgenetic alopecia affects the majority of men to some degree during their lifetime. The prevalence increases with age: approximately 20\% of men in their 20s, 30\% in their 30s, and up to 50\% or more of men over the age of 50 experience noticeable male-pattern baldness. While it is a global phenomenon, the prevalence and severity can vary among different ethnic groups, with Caucasian men generally having higher rates of hair loss compared to Asian and Afro-Caribbean men.

\section*{Aetiology and Pathophysiology}
Male-pattern baldness is caused by an inherited sensitivity of scalp hair follicles to male hormones:

\begin{itemize}
    \item \textbf{Genetic predisposition:} Polygenic inheritance from both maternal and paternal lineages determines the sensitivity of hair follicles to androgens.
    \item \textbf{Dihydrotestosterone (DHT):} The enzyme 5-alpha reductase converts testosterone into DHT, which binds to androgen receptors in susceptible hair follicles.
    \item \textbf{Follicular miniaturisation:} Under the influence of DHT, the active growth phase (anagen) of the hair cycle shortens, and the hair follicles shrink, producing progressively finer, shorter, and less visible hairs.
    \item \textbf{Terminal to vellus transition:} Eventually, thick pigmented terminal hairs are replaced by tiny, non-pigmented vellus hairs, and the follicles may eventually stop producing hair altogether.
\end{itemize}

\section*{Clinical Features}
Male-pattern baldness progresses through well-defined stages (often classified using the Norwood scale) with typical presentations:

\begin{itemize}
    \item \textbf{Receding hairline:} Thinning and recession beginning at the temples, creating an "M" shape when viewed from the front.
    \item \textbf{Vertex thinning:} Thinning or a bald patch appearing at the crown of the head.
    \item \textbf{Progressive confluence:} The receding hairline and bald crown expand and eventually join, leaving hair only on the sides and back of the scalp.
    \item \textbf{Normal scalp skin:} The underlying scalp skin remains healthy, without redness, scaling, pain, or scarring.
\end{itemize}

\section*{Differential Diagnosis}
Other types of hair loss must be differentiated from androgenetic alopecia:

\begin{itemize}
    \item \textbf{Alopecia areata:} Autoimmune hair loss presenting as sudden, patchy, circular bald spots, which can occur anywhere on the body.
    \item \textbf{Telogen effluvium:} Temporary, diffuse hair shedding across the entire scalp, often triggered by severe physiological stress, illness, or medications.
    \item \textbf{Tinea capitis:} A fungal infection of the scalp causing patchy hair loss with scale, redness, and broken hairs.
    \item \textbf{Traction alopecia:} Hair loss caused by chronic pulling or tension on the hair shafts due to tight hairstyles.
    \item \textbf{Anabolic steroid use:} Exogenous androgens accelerating hair loss in genetically predisposed individuals.
\end{itemize}

\section*{When to Seek Medical Care}
Men concerned about progressive hair loss or those wishing to discuss treatment options should consult a GP or dermatologist. Early intervention is key, as treatments are more effective at preserving existing hair than regrowing lost hair. Immediate medical care is not required for male-pattern baldness itself, but sudden or atypical hair loss requires clinical investigation.

\textbf{Call 000 (or local emergency services) for:}
\begin{itemize}
    \item Severe, sudden facial or throat swelling accompanying a new medication for hair loss (allergic reaction)
    \item Difficulty breathing or chest pain after starting a new oral medication
    \item Sudden loss of consciousness or severe dizziness
    \item A high fever accompanied by a rapidly spreading, painful scalp infection
    \item Any acute, life-threatening emergency symptom
\end{itemize}

\section*{Diagnosis and Investigations}
Diagnosis is primarily clinical, based on the pattern of hair loss and clinical history:

\begin{itemize}
    \item \textbf{Scalp examination:} Evaluating the pattern of hair loss and examining the hair shafts and scalp skin using a dermatoscope (trichoscopy).
    \item \textbf{Medical history:} Reviewing the speed of hair loss, family history of baldness, diet, medications, and systemic symptoms.
    \item \textbf{Pull test:} Gently pulling a small bundle of hair to assess the active shedding rate and rule out telogen effluvium.
    \item \textbf{Blood tests:} Checking thyroid function, iron studies, or hormone levels if the pattern is atypical or other systemic illness is suspected.
\end{itemize}

\section*{Management}
Several options exist to manage, slow down, or cosmetically treat male-pattern baldness:

\begin{itemize}
    \item \textbf{Topical minoxidil:} An over-the-counter liquid or foam applied directly to the scalp twice daily to stimulate hair growth and prolong the growth phase.
    \item \textbf{Oral finasteride:} A prescription tablet that inhibits 5-alpha reductase, reducing scalp DHT levels by about 60--70\%, helping to stop hair loss and promote regrowth.
    \item \textbf{Hair transplant surgery:} A surgical procedure (FUT or FUE) where hair follicles are transferred from the back and sides of the scalp to the thinning areas.
    \item \textbf{Low-level laser therapy (LLLT):} Using red light devices to stimulate cellular activity in hair follicles.
    \item \textbf{Cosmetic solutions:} Wigs, hairpieces, scalp micro-pigmentation (tattooing), or styling techniques to conceal thinning areas.
\end{itemize}

\section*{Complications}
\begin{itemize}
    \item Psychological distress, including reduced self-esteem, anxiety, depression, and poor body image
    \item Increased risk of sunburn and skin cancers on the exposed, unprotected scalp skin
    \item Potential side effects of finasteride, including sexual dysfunction (decreased libido or erectile dysfunction) in a small percentage of users
    \item Side effects of topical minoxidil, such as scalp irritation, itching, or unwanted facial hair growth if misapplied
    \item Surgical risks of hair transplantation, including infection, scarring, and unnatural-looking results
\end{itemize}

\section*{Prognosis and Outcomes}
Male-pattern baldness is a chronic, progressive condition that continues throughout life if untreated. Treatments like minoxidil and finasteride can successfully halt hair loss and improve hair density in many men, but they must be used continuously; stopping treatment leads to a return of the hair loss process within 6 to 12 months. Hair transplant surgery provides permanent results for the transplanted follicles, though surrounding non-transplanted hair may continue to thin.

\section*{Prevention and Self-Care}
\begin{itemize}
    \item Protect the bald areas of the scalp from ultraviolet radiation by wearing a hat and applying SPF 50+ sunscreen
    \item Wash the scalp regularly with a gentle, non-irritating shampoo to maintain skin health
    \item Avoid harsh chemical treatments, excessive heat styling, or vigorous brushing, which can break fragile hairs
    \item Eat a balanced diet rich in essential vitamins, minerals, and proteins to support overall hair health
    \item Avoid smoking, which may impair microcirculation to hair follicles and accelerate hair thinning
\end{itemize}

\section*{Sources and Further Reading}
The following organisations provide reliable information on male-pattern baldness:

\begin{itemize}
    \item Australasian College of Dermatologists. \textit{A-Z guide on male androgenetic alopecia and treatments.} https://www.dermcoll.edu.au/
    \item NHS. \textit{Patient information on hair loss causes and treatments.} https://www.nhs.uk/
    \item Mayo Clinic. \textit{Guide to hair loss, diagnosis, and medical management.} https://www.mayoclinic.org/
    \item Healthdirect Australia. \textit{Male pattern baldness symptoms, treatments, and emotional support.} https://www.healthdirect.gov.au/
    \item British Association of Dermatologists. \textit{Patient information leaflet on androgenetic alopecia.} https://www.bad.org.uk/
\end{itemize}
